\chapter{Materiais e métodos}
\label{capitulo:materiais_metodos}

Os materiais e métodos compõem a metodologia do trabalho.
Neste momento deverão ser apresentadas as principais tecnologias utilizadas.

Entende-se como materiais: 

\begin{itemize}
    \item linguagens de programação (Julia, Java, PHP, Python);

    \item \textit{frameworks} ou arcabouços;

    \item Banco de dados (PostgreSQL, MySQL, SQL Server, Oracle, etc.);

    \item Ambiente de Desenvolvimento Integrado (IDE): NetBeans, Eclipse, etc;

    \item APIs de sistemas;

    \item \textit{Hardwares}: seu computador\footnote{Sim! Seu computador. Se estiver fazendo um trabalho que envolva calcular tempo de execução de um algoritmo, as definições do seu computador devem ser incluídas no trabalho, tais como: processador, número de núcleos, número de \textit{threads}, memória RAM, tipo de disco rígido (HD ou SSD), tamanho do disco rígido, até mesmo a velocidade de conexão deve ser incluída dependendo do trabalho.}, Arduíno, Raspberry Pi, etc.
\end{itemize}
    
Não entenda mal, a questão deste capítulo não consiste apenas em citar os materiais, mas sim falar onde e como eles foram utilizados na sua pesquisa.
Então ao citar o computador, deve ser mencionado nos métodos como o computador foi utilizado.

Como métodos entende-se:

\begin{itemize}
    \item metodologia de desenvolvimento (SCRUM, cascata, incremental, etc.);
    
    \item qual forma o sistema foi modelado? Usou-se Diagrama de Fluxo de Dados (DFD)? Unified Modeling Language (UML)? O que foi usado?

    \item Como os requisitos foram levantados? Quais técnicas foram usadas nas entrevistas? Onde foram as entrevistas? Quando? Quem eram os envolvidos?

    \item Quais etapas foram seguidas no desenvolvimento?
\end{itemize}

O autor deve tomar cuidado para não confundir os materiais com a revisão bibliográfica (capítulo \ref{capitulo:revisao_bibliografica}).
Note que neste capítulo não é o momento de definir o que são os materiais, mas sim como eles foram utilizados.
Então se o autor vai apresentar a linguagem Julia ao leitor, isto deve ser feito no capítulo de revisão bibliográfica na parte de fundamentação teórica (seção \ref{secao:fundamentacao_teorica}).